
Predicting properties of trained neural networks directly from their weights --- generalization gap, test accuracy, training loss --- has applications in model selection, hyperparameter analysis, and automated machine learning. Unterthiner et al. (2020) demonstrated that a flat multi-layer perceptron (flat-MLP) applied to concatenated weight vectors achieves high predictive performance on large model zoos, establishing this as the standard baseline for weight-space property prediction.

A structural property of fully-connected neural networks is that any permutation of neurons within a layer yields a functionally identical network. Flat-MLP encoders, which concatenate weights into a fixed-order vector, treat neuron position as an input feature. When the neuron ordering is consistent across models (as in a single training pipeline), this does not prevent effective prediction. However, when neuron orderings differ --- due to different random seeds, training procedures, or deliberate permutation --- flat-MLP representations may change for functionally equivalent networks.

Neural Functional Transformers (NFT; Zhou et al., 2023) address this by applying permutation-equivariant multi-head attention over per-neuron token sequences. The equivariance property guarantees that any permutation-invariant downstream head produces identical predictions regardless of neuron ordering. NFT was previously evaluated on implicit neural representation (INR) classification tasks but had not been tested on model zoo property prediction tasks such as generalization gap regression.

This work bridges these two lines of research by conducting a controlled comparison of NFT and flat-MLP encoders for generalization gap prediction on the Unterthiner model zoo. The specific contributions are:

\begin{enumerate}
\item A direct empirical comparison of NFT and flat-MLP encoders under permutation stress for generalization gap prediction, using a six-encoder ablation suite spanning no equivariance handling through oracle canonicalization.
\end{enumerate}

\begin{enumerate}
\item A mediation analysis (following Baron and Kenny, 1986) testing whether the observed robustness difference is attributable to equivariant attention capturing neuron influence concentration signals, yielding delta-R-squared = 0.228.
\end{enumerate}

\begin{enumerate}
\item Empirical evaluation of two alternative approaches: permutation augmentation (partial reduction with high variance) and L2-norm canonicalization (collapsed to constant predictor).
\end{enumerate}

\begin{enumerate}
\item An observation that NFT achieved higher baseline Spearman rho (0.489 vs. 0.303) with 40x fewer parameters (75K vs. 3.04M), though this comparison is confounded by unmatched parameter counts (see Section 6).
\end{enumerate}

The scope of this study is limited to within-distribution evaluation on a single model zoo. Planned cross-pipeline transfer experiments (MNIST to CIFAR) were not executed and are noted as future work.
